\section{实验设计}
\label{zh:sec:experimental_design}

实验评估在三个不同层次展开：大规模开环统计仿真、闭环滚动重规划以及在 Apollo Planning 生产架构中的软件集成验证。统计实验使用确定性场景生成器和固定随机种子，使每种方法接收完全相同的环境输入，碰撞与间隙指标则依据独立保留的真实轨迹进行严格评估。

\subsection{实验协议与场景}
\label{zh:sec:protocols_scenarios}

每个场景使用按场景族分配的固定时域 $T\in\{9,10,12,14\}$ s：交叉行人与延迟交叉为 9 s，车辆切入、带对向交通的停放车辆及预测噪声为 10 s，窄间隙多障碍物为 12 s，交错动态障碍物为 14 s。当仿真时钟达到 $T$ 而未达到 $s_{\mathrm{end}}\geq s_{\max}-1$ m 时判为超时。全文使用四套互不重叠的协议。

\emph{开环评估协议，涵盖 3{,}500 个场景。} 七类场景族各使用 500 个随机种子：交叉行人、车辆切入、带对向交通的停放车辆、窄间隙多障碍物、交错动态障碍物、预测噪声与延迟交叉。参数在各场景族范围内采样。除预测噪声场景族仅将有界误差注入规划输入外，规划场景与真实场景完全一致。

\emph{固定截面几何检查，涵盖 4{,}000 个实例。} 20 个确定性随机种子，每个种子生成 200 个固定截面区间实例；每个实例包含 0 到 8 个可能重叠、嵌套或穿越道路边界的禁行区间。扫描算法与全部 $2^k$ 种侧向分配的穷举结果比较中心带宽度。

\emph{滚动重规划评估协议，涵盖 700 个场景。} 每类场景族使用 100 个随机种子。自车在 $t=0$ 时由真实状态 $(s_0,l_0,v_0,a_0)$ 初始化，首个规划周期采用冷启动，不复用 QP 对偶变量的暖启动，也不承继任何同伦决策。每隔 0.5 s 执行当前计划的第一个片段，更新障碍物状态，并在移位后的时域 $\max(2.0\text{ s},T-t)$ 上再次调用规划器，直至自车在 $s\geq s_{\max}-1$ m 到达目标或达到 $T$。

\emph{有限网格联合参考，涵盖 70 个场景。} 分层设置，每类场景族使用 10 个随机种子，网格间距为 $\Delta t=0.5$ s、$\Delta s=0.5$ m、$\Delta l=0.25$ m 与 $\Delta v=1.0$ m/s，束宽设为 20{,}000。

\subsection{基线与实现}
\label{zh:sec:baselines_impl}

匹配的 PVD 基线按照约束接口交换的信息定义：
\begin{itemize}
  \item B0 无时间意识，使用固定的局部侧向策略。
  \item B1 增加依赖到达时间的投影，但保留逐障碍物的贪心侧向选择。
  \item B2 在 B1 策略中允许三次到达时间细化迭代。
  \item MIKU 增加群组级连续划分、Top-$K$ 空间同伦、威胁感知裕度、不确定性扩展占用、时间同伦，以及最多两次交替更新的规划迭代。
\end{itemize}
四种方法共享相同的路径 QP 与速度 QP 目标、运动学约束、输入及指标实现。两个分段加加速度 QP 均在 OSQP 的默认容差下求解，3{,}500 场景运行中不施加逐次调用的墙钟上限。基线 B2 在其三次到达时间细化迭代内使用同一套 QP。基线 B3 是 $(t,s,l,v)$ 网格上的确定性集束搜索动态规划，束宽设为 20{,}000。单步转移代价为相对 $v_0$ 的速度偏差、纵向加速度、横向偏置与横向加速度的加权和再减去进度奖励；运动学不可行或与不确定性扩展障碍物发生碰撞的转移在排序前直接剪除。MIKU 的全部超参数在评估前固定，并汇总于表~\ref{zh:tab:fixed_parameters}。实验运行于配有 64 GB 内存的 Intel i9-14900HX 平台。Apollo 集成验证与 Apollo 11.0 Planning 的兼容性。

\begin{table}[t]
\centering
\caption{全部实验采用的固定规划与预测误差参数}
\label{zh:tab:fixed_parameters}
\small
\begin{tabular}{lll}
\toprule
参数 & 数值 & 用途\\
\midrule
$(w_1,\ldots,w_5)$ & $(0.30,0.20,0.15,0.10,0.25)$ & 威胁度权重\\
TTC 断点 & $2.0,7.0$ s & 紧迫度归一化\\
$(\delta_{\min},\delta_{\max})$ & $(0.10,0.40)$ m & 自适应障碍物缓冲\\
$Q,K_s,\lambda_s$ & $3,3,0.05$ & 空间带、Top-$K$ 与切换权重\\
$B$ & $8$ & 时间图束宽\\
$\eta_c$ & $0.10+0.05\min(|v_l|,2)$ s & 占用时间窗扩展\\
$\epsilon_s$ & $0.05$ m & 纵向 ST 安全裕度\\
$(\epsilon_{s0},\epsilon_{l0})$ & $(0.35,0.18)$ m & 预测位置初始误差界\\
$(\epsilon_{vs},\epsilon_{vl})$ & $(0.30,0.12)$ m/s & 预测误差增长率界\\
路径/速度离散步长 & $0.5$ m / $0.1$ s & QP 采样\\
规划迭代次数 & 最多 $2$ 次 & 初始求解与一次细化\\
\bottomrule
\end{tabular}
\end{table}

\subsection{结果指标与统计分析}
\label{zh:sec:outcomes_stats}

主要结果指标是无碰撞到达目标率：当运行在完整时域内保持无碰撞，并且达到 $s_{\mathrm{end}}\geq s_{\max}-1$ m 时判为成功。碰撞率根据真实矩形评估；退化停车率记录无碰撞但停车或未到达目标的运行。进度比、平均速度、加加速度均方根、最小带符号间隙、最小纵向碰撞时间、带惩罚的行驶时间与规划延迟提供安全性、舒适性与计算方面的背景。无碰撞到达时，带惩罚的行驶时间等于到达时间，否则等于第~\ref{zh:sec:protocols_scenarios} 节中各场景族对应的时域 $T$ 加 5 s。运行时间涵盖完整的方法调用，B2 包括其全部协调迭代。

对于每个固定随机种子，先计算 MIKU 减去基线的差值再行聚合。连续指标使用 5{,}000 次重采样的配对百分位 Bootstrap 区间；二元结果使用精确 McNemar 检验；效应量采用配对 Cohen's $d_z$。运行时间报告第 50、95 与 99 百分位。
